\subsection{Equal finite metrics and asymptotic speed with different limiting moduli}
\label{sec:same-speed-different-modulus}

The exact rank-two pair also separates asymptotic speed
from marked metric scale. Both branches have the same
limiting speed function. Nevertheless a suitable two-mode
phase gives different circle means of the limiting
metric potential even when both induced metrics
were equal at the measured times.

Let $w(\tau,\theta)=A(\theta)\tau+B(\theta)+o(1)$ be
a finite sum of Dirichlet-node seed modes, and let
$\sigma$ be the common speed of the two seed geodesics
with orthonormal initial tangent components $(\pm g,q)$.
Assume $g>0$ and $q\ne0$.
The seed metric potential is identical on both histories.
Write
\[
 D=\log\frac{1+g/\sigma}{1-g/\sigma}.
\]
The explicit hyperbolic geodesics give, wherever
$A(\theta)\ne0$,
\[
 \lim_{\tau\to\infty}
     (\widetilde P_+-\widetilde P_-)
      =-D\operatorname{sgn} A(\theta),\qquad
 \lim_{\tau\to\infty}
     (\widetilde\lambda_+-\widetilde\lambda_-)
      =-2D\operatorname{sgn} A(\theta).
\]
Indeed $e^{P_\pm-P_0}
=\cosh(\sigma w)\pm(g/\sigma)\sinh(\sigma w)$,
and the Ernst transformation changes the metric
potential difference by the same factor two.
The transformed limiting speed of each branch is
$1-\sigma|A(\theta)|$ for sufficiently small amplitude.
At a zero of $A$ the position difference instead tends
to its finite value evaluated at $w=B(\theta)$.
At a simple zero $\theta_0$, the limit with
$\theta=\theta_0+\xi/\tau$ is the corresponding
explicit geodesic ratio evaluated at
$B(\theta_0)+A'(\theta_0)\xi$.
This gives the angular scale of the marked profile
transition. Wherever $A\ne0$, the ratio of the marked
$x$-cycle lengths tends to
$\exp(-D\operatorname{sgn} A/2)$.

\begin{proposition}
\label{prop:same-speed-different-modulus}
There are smooth vacuum three-torus Gowdy pairs with
rank-two target differential on open sets, identical
induced metrics and $Q_t$ at the two measured slices,
and identical asymptotic speed at every angle, whose
limiting metric potentials have different circle means.
\end{proposition}

\begin{proof}
Choose a radial annulus ratio $R$ such that the fifth
Dirichlet frequency is six times the first.
For the rescaled outer radius one, let $z_j=R\nu_j(R)$.
Interval brackets for the first and fifth roots of
$J_0(z/R)Y_0(z)-Y_0(z/R)J_0(z)$ are
\[
\begin{array}{c|cc}
 R & z_1 & z_5\\ \hline
64 & (2.850,2.852) & (15.737,15.739)\\
10^{10} & (2.473,2.475) & (15.007,15.009).
\end{array}
\]
These are verified by the power series for $J_0,Y_0$
with interval arithmetic and geometric bounds on the
series tails. The indices follow from
$j_{0,j}\le z_j\le j\pi/(1-1/R)$:
the lower bound is comparison with the disk by zero
extension in the hole and the upper bound follows
from the transformed radial operator
$-\partial_r^2-1/(4r^2)$.
Continuity gives the required ratio at some $R>64$.
Set $a=z_1/R$ and $b=z_1$, so the angular modes
$1$ and $6$ have common nodes. The upper bound on $z_1$
gives $6a<0.30$, and both relevant $J_0$ values at the
inner radius are positive.

We can then use the two scalar modes with coefficients
chosen to give
\[
 A(\theta)=A_0\big(\cos\theta+\eta\cos(6\theta+\phi)\big)
\]
for small $A_0>0$, small nonzero $\eta$, and
$\phi/\pi$ irrational with $\cos\phi\ne0$.
Indeed the logarithmic slope of each mode is a nonzero
multiple of $J_0(ka)$. The sign-positive set has
length $\pi-2\eta\cos\phi+O(\eta^2)$ by
implicit differentiation at the two zeros
$\pi/2$ and $3\pi/2$. Thus
$\overline{\operatorname{sgn}A}\ne0$.
The initial and terminal metrics and velocities
obey Theorem~\ref{th:rank-two}; the same rank
argument works for the two-mode phase.

To pass to the circle mean, the finite-Bessel seed
expansion is uniform and its metric remainder is bounded.
Also $P_\pm(w)-\sigma|w|$ is uniformly bounded,
so the transformed potentials after subtracting their
common leading $\tau(1-\sigma|A|)^2$ have pointwise
limits and uniform domination. Their circle-mean
difference is therefore
$-2D\,\overline{\operatorname{sgn}A}\ne0$.
\end{proof}

The ratio argument concerns the marked regular-slice
geometry and the displayed asymptotic observables.
It does not infer a change of global spacetime fate
from the velocity-one sites or from the coordinate scale
of the transition.
